\section{Methodology}

\subsection{Problem Formulation}

Let $x \in \mathbb{R}^{C \times T}$ denote a multi-lead ECG recording with
$C$ leads and $T$ temporal samples ($C$=12, $T$=1{,}000 at 100\,Hz). The
multi-label task assigns $x$ a target $y \in \{0,1\}^K$ over $K$=5
superdiagnostic classes. An encoder $F(\cdot\,;\theta_F)$ maps the input to an
embedding $z = F(x;\theta_F) \in \mathbb{R}^d$.

Conventional practice fixes $d$ and trains $N$ separate encoders
$\{F_1,\dots,F_N\}$ for $N$ deployment tiers, incurring $N\cdot|\theta_F|$
parameters and $N$ independent validation workflows. The nested alternative
trains one encoder whose prefixes are individually usable.

\subsection{Signal Preprocessing}
\label{sec:preproc}

Let $x^{\mathrm{raw}} \in \mathbb{R}^{C\times T}$ be the recording in physical
units (mV). We standardise each lead using statistics computed once over the
\emph{training split}:
\begin{equation}
x^{\mathrm{norm}}_{c,t} = \frac{x^{\mathrm{raw}}_{c,t} - \mu_c}{\sigma_c},
\qquad
\mu_c, \sigma_c \ \text{from} \ \mathcal{D}_{\mathrm{train}}.
\label{eq:norm}
\end{equation}

This differs from the per-record, per-lead $z$-scoring used in the earlier
version of this work, and the difference is consequential. Per-record scaling
divides each lead by its own standard deviation, which removes inter-patient
amplitude variation but also \emph{destroys absolute voltage}. Voltage
amplitude is the defining criterion for ventricular hypertrophy---the
Sokolow--Lyon and Cornell indices are threshold tests on millivolt
sums---so per-record normalisation discards precisely the evidence the HYP
class depends on. We report this as an ablation
(Section~\ref{sec:ablations}) because it plausibly explains both the weak HYP
performance and part of the gap to the published PTB-XL benchmark.

Training-time augmentation applies additive Gaussian noise
($\sigma_n$=0.01, $p$=0.5), amplitude scaling ($\alpha \sim
\mathcal{U}(0.95,1.05)$, $p$=0.5), circular temporal shift ($\delta \sim
\mathcal{U}(-10,10)$ samples, $p$=0.3), and lead dropout zeroing one or two
leads ($p$=0.1). No bandpass filtering is applied: at 100\,Hz the Nyquist
limit already constrains the bandwidth to 50\,Hz.

\subsection{Matryoshka Objective}

Define nesting dimensions $\mathcal{M} = \{m_1,\dots,m_{|\mathcal{M}|}\}$ with
$m_1 < \cdots < m_{|\mathcal{M}|} = d$. For each $m$, let $g_m:\mathbb{R}^m
\to \mathbb{R}^K$ be a classifier operating on the prefix $z_{1:m}$. The
objective is
\begin{equation}
\ell_{\mathrm{MRL}}(\theta_F, \{\theta_g\}) =
\frac{1}{|\mathcal{M}|}\sum_{m \in \mathcal{M}}
w_m \cdot \ell_{\mathrm{BCE}}\big(g_m(z_{1:m}),\, y\big),
\label{eq:mrl}
\end{equation}
with $w_m$ = 1 uniformly and $\mathcal{M} = \{16,32,64,128,256,512\}$
throughout. Alternative weightings are ablated in
Section~\ref{sec:ablations}.

Because $z_{1:m}$ is a strict prefix of $z_{1:m'}$ for $m<m'$, coordinate $j$
receives gradient from every head with $m \ge j$:
\begin{equation}
\mathbb{E}\!\left[\left\|\frac{\partial \ell_{\mathrm{MRL}}}{\partial z_j}\right\|\right]
= \frac{1}{|\mathcal{M}|}\!\!\sum_{\substack{m \in \mathcal{M} \\ m \ge j}}\!\!
\mathbb{E}\!\left[\left\|\frac{\partial \ell_{\mathrm{BCE}}(g_m(z_{1:m}),y)}{\partial z_j}\right\|\right]\!.
\label{eq:gradpressure}
\end{equation}

Since $|\{m \in \mathcal{M}: m \ge j\}|$ decreases monotonically in $j$,
leading coordinates receive strictly more gradient signal. This is a
\emph{structural} property of the objective, and it motivates the hypothesis
that leading coordinates encode the most broadly discriminative features. It
does not, however, establish that they encode any \emph{particular} clinical
feature---a stronger claim requiring measurement, which
Section~\ref{sec:probing} supplies.

We use MRL with independent per-dimension heads, and additionally evaluate the
efficient variant MRL-E, which replaces $|\mathcal{M}|$ classifiers with a
single matrix $W \in \mathbb{R}^{K \times d}$ column-sliced per granularity:
\begin{equation}
g^{(E)}_m(z_{1:m}) = W_{:,1:m}\, z_{1:m}.
\label{eq:mrle}
\end{equation}
This reduces classifier parameters from $K\sum_m m + K|\mathcal{M}|$ to
$K d + K|\mathcal{M}|$.

\subsection{Backbone Architectures}

\textbf{Inception1D}~\cite{ismailfawaz2020} uses parallel convolutions with
kernel sizes 9, 19 and 39 plus a max-pooled branch, concatenated
channel-wise after a $1\times1$ bottleneck:
\begin{equation}
h_{\mathrm{inc}} = \mathrm{BN}\big(\mathrm{Cat}[\mathrm{Conv}_9(b),
\mathrm{Conv}_{19}(b), \mathrm{Conv}_{39}(b), \mathrm{MaxPool}(x)]\big),
\end{equation}
with $b = \mathrm{Conv}_{1\times1}(x)$. Squeeze-and-excitation~\cite{hu2018}
is applied within each block. We note a correction: the earlier version of
this work described SE attention in the Inception1D configuration but the
released implementation did not include it. SE is now genuinely enabled and
the parameter counts reported here reflect that.

\textbf{XResNet1D-101}~\cite{he2016,strodthoff2021} uses a three-layer
convolutional stem (kernels 7, 3, 3) followed by four residual stages of
$[3,4,23,3]$ SE-bottleneck blocks, global average pooling and a linear
projection to $d$.

Both terminate in a $d$-dimensional embedding followed by batch
normalisation. Exact parameter counts, measured rather than quoted, appear in
Table~\ref{tab:compute}.

\subsection{Training Procedure}

Algorithm~\ref{alg:train} summarises training. All heads are optimised jointly
in a single backward pass. We use AdamW~\cite{loshchilov2019} with cosine
annealing after linear warmup, gradient clipping at norm 1.0, and weight decay
applied only to parameters with $\mathrm{ndim} > 1$ (a correction: the
previous implementation's substring-based parameter grouping applied decay to
normalisation parameters nested inside convolutional blocks).

Model selection uses the \emph{mean} validation macro AUC across all
granularities rather than the AUC at $d$=512. Selecting on the largest head
alone biases the checkpoint toward one operating point, which is exactly the
behaviour a nested model should avoid.

\begin{algorithm}[!t]
\caption{MRL-ECG training}
\label{alg:train}
\begin{algorithmic}[1]
\Require training set $\mathcal{D}$, nesting dims $\mathcal{M}$, epochs $E$,
patience $P$
\Ensure $\theta_F^*$, $\{\theta_{g_m}^*\}_{m\in\mathcal{M}}$
\State initialise encoder $F$ and heads $\{g_m\}$; $a^* \gets 0$; $w \gets 0$
\For{$e = 1$ to $E$}
  \For{each mini-batch $\mathcal{B} \subset \mathcal{D}$}
    \State $\mathcal{B}' \gets \mathrm{Augment}(\mathrm{Normalise}(\mathcal{B}))$
    \State $Z \gets F(\mathcal{B}'; \theta_F)$
    \State $\ell \gets \frac{1}{|\mathcal{M}|}\sum_{m} w_m
           \ell_{\mathrm{BCE}}(g_m(Z_{:,1:m}), y)$
    \State update $\theta_F, \{\theta_{g_m}\}$ via AdamW on $\nabla\ell$
  \EndFor
  \State $a \gets \frac{1}{|\mathcal{M}|}\sum_m
         \mathrm{macroAUC}(g_m, \mathcal{D}_{\mathrm{val}})$
         \Comment{mean over granularities}
  \If{$a > a^*$}
    \State $a^* \gets a$; checkpoint; $w \gets 0$
  \Else
    \State $w \gets w+1$; \textbf{if} $w \ge P$ \textbf{then break}
  \EndIf
\EndFor
\State tune per-class thresholds on $\mathcal{D}_{\mathrm{val}}$; freeze for test
\end{algorithmic}
\end{algorithm}

Decision thresholds for F1 are tuned per class on the validation split and
then frozen. The earlier version evaluated F1 at a fixed threshold of 0.5,
which for an imbalanced multi-label problem understates F1 substantially and
accounts for most of the apparent discrepancy between the reported AUC
($\approx$0.90) and F1 ($\approx$0.62).

\subsection{Deployment Mapping}
\label{sec:deployment}

Given a device embedding budget $B_\tau$ bytes, the admissible truncation
point is
\begin{equation}
\pi(\tau) = \max\{m \in \mathcal{M} : 4m \le B_\tau\},
\label{eq:deploy}
\end{equation}
the factor 4 accounting for float32 storage. We present
Equation~\ref{eq:deploy} as arithmetic, not as a validated deployment
result. It states which truncation point \emph{fits} a given budget; whether
the resulting diagnostic behaviour is acceptable on data acquired by that
device is a separate empirical question, addressed only partially by the
reduced-lead and artefact experiments of Section~\ref{sec:wearable} and not at
all by execution on microcontroller hardware, which we have not performed.
